\section{Requirements}

\subsection{Requirement Mapping Strategy}
The assignment statement provides high-level tasks, but a report must explain
how those tasks become concrete, testable requirements in the implemented
system. In this project, each assignment task is refined into two layers:
\begin{enumerate}
  \item a functional requirement that can be observed by the user, and
  \item an implementation-oriented requirement that constrains timing, task
  communication, or system architecture.
\end{enumerate}
This section follows the same logic. It first describes the detailed
requirements of each task and then summarizes the non-functional properties that
the complete system must satisfy.

\subsection{Task 1 Requirements: Temperature-Dependent LED Behavior}
The assignment requires at least three distinct LED behaviors based on
temperature. In the implemented system, the following concrete requirements were
derived:
\begin{itemize}
  \item The system shall classify temperature into at least three operating
  regions.
  \item The discrete LED shall blink with a different cadence for each
  temperature region.
  \item LED mode changes shall be triggered automatically whenever a new valid
  temperature sample is produced.
  \item The LED logic shall run independently of the sensor task so that the
  sensing period remains stable.
  \item The LED task shall also support short event pulses for system feedback
  such as Wi-Fi-connected or error-notification events.
\end{itemize}

The current implementation satisfies these requirements by mapping temperature
to three primary modes: cold, normal, and hot. These modes correspond to
significantly different blink half-periods, which makes the environmental state
easy to distinguish even without reading the LCD or web dashboard.

\subsection{Task 2 Requirements: Humidity-Dependent NeoPixel Indication}
For the NeoPixel task, the assignment requires humidity-based color indication
using semaphore-style synchronization. In the implemented project, the
requirement was refined as follows:
\begin{itemize}
  \item The system shall compute a humidity-dependent color output for the
  NeoPixel.
  \item The color shall reflect at least three semantic humidity regions, even
  if the implementation internally uses a continuous gradient.
  \item The NeoPixel update shall be decoupled from sensing through a queue so
  that the LED strip always displays the most recent humidity state.
  \item The NeoPixel task shall support an enable/disable command and clear the
  LED when disabled.
\end{itemize}

In the current code, humidity values are constrained to the range
\(0\%\) to \(100\%\) and converted into a color gradient. The lower half of the
range is represented by a transition from red to green, while the upper half is
represented by a transition from green to blue. This means the assignment's
``at least three levels'' condition is exceeded by a smooth visual scale.

\subsection{Task 3 Requirements: LCD-Based Monitoring and RTOS
Synchronization}
The LCD task was extended beyond a static two-line text display. The refined
requirements are:
\begin{itemize}
  \item The LCD shall display valid temperature and humidity measurements when
  the sensor is operating normally.
  \item The LCD shall display an error message when sensing fails.
  \item The LCD shall support more than one display state, including a sensor
  mode and a network-aware mode.
  \item In STA mode, the LCD shall be able to switch between raw sensor data
  and TinyML forecast information.
  \item The DHT20 sensor and LCD share the same I\textsuperscript{2}C bus, so
  access to that bus shall be protected by a mutex.
\end{itemize}

The implemented design therefore uses the LCD not only as an output device but
also as a demonstration of synchronization. The I\textsuperscript{2}C mutex is
one of the clearest examples in the code where a shared hardware resource is
explicitly protected.

\subsection{Task 4 Requirements: AP-Mode Web Server and User Interaction}
The web server task is one of the most visible parts of the project and also
one of the most functionally rich. The concrete requirements implemented in the
current version are:
\begin{itemize}
  \item The board shall provide an HTTP server when running in AP mode.
  \item The root dashboard shall show temperature, humidity, TinyML prediction,
  Wi-Fi state, and chart history.
  \item The web server shall provide Wi-Fi scanning so the user can list
  available networks.
  \item The web server shall provide storage and recall of saved Wi-Fi
  credentials.
  \item The interface shall provide at least two direct control actions. In
  this project those actions are:
  \begin{itemize}
    \item external LED on/off toggle,
    \item external LED brightness adjustment through a PWM slider.
  \end{itemize}
  \item The system shall be able to move from AP mode to STA mode using
  browser-supplied SSID and password.
  \item If STA connection does not succeed within the timeout window, the
  system shall return to AP mode automatically.
  \item Pressing the BOOT button shall force the board back to AP mode.
\end{itemize}

The web task therefore covers both UI and network-state control. This is more
than a simple monitoring page: it acts as the configuration front-end for the
entire device.

\subsection{Task 5 Requirements: TinyML Deployment and Evaluation}
The TinyML portion of the assignment requires a complete mini-pipeline from
data representation to inference. The concrete requirements implemented in code
are:
\begin{itemize}
  \item The system shall deploy a TensorFlow Lite Micro model on the ESP32-S3.
  \item The TinyML task shall receive live sensor data from the RTOS pipeline.
  \item Input features shall be transformed and normalized before inference.
  \item The model output shall be converted into an interpretable prediction.
  \item The prediction result shall be accessible to other tasks, especially
  the LCD task, the web task, and the CoreIOT task.
  \item The report shall discuss the expected inference timing and the impact
  of batching and feature engineering on accuracy and responsiveness.
\end{itemize}

The implemented TinyML logic uses batch averaging and derived time-based
features rather than passing raw sensor samples directly to the model. This is
important because it shows that TinyML integration is treated as a system
design problem, not merely as a call to an inference library
\cite{tflm}.

\subsection{Task 6 Requirements: CoreIOT Cloud Publishing}
The cloud task is constrained by the assignment's STA-mode rule. The refined
requirements are:
\begin{itemize}
  \item The board shall connect to CoreIOT only when Wi-Fi STA mode is active
  and Internet access is available.
  \item The device shall use a valid CoreIOT token to authenticate.
  \item Telemetry shall include temperature and humidity, and may include
  additional derived fields that improve observability.
  \item The cloud task shall periodically publish data without blocking local
  sensing and control.
  \item The system shall subscribe to RPC messages so that a remote cloud
  command can influence local behavior.
\end{itemize}

In the current implementation, the RPC channel is used to control whether the
LCD should show the TinyML forecast in STA mode. This is a strong design choice
because it demonstrates bidirectional interaction with the cloud rather than
simple telemetry upload \cite{coreiot_platform,pubsubclient,arduinojson}.

\subsection{Cross-Cutting Functional Requirements}
Beyond the six assignment tasks, the following system-level functional
requirements emerge from the integrated design:
\begin{itemize}
  \item The latest sensor sample shall be distributed to multiple consumers
  without unnecessary duplication of acquisition logic.
  \item The user shall be able to understand the state of the system through
  local logs, LCD messages, and browser-visible status labels.
  \item The system shall preserve the most recent Wi-Fi profiles across resets
  by using non-volatile storage.
  \item The dashboard shall plot recent sensor history rather than showing only
  the most recent sample.
  \item The cloud payload shall include both raw values and interpreted state
  such as rain probability and LCD forecast mode.
\end{itemize}

\subsection{Non-Functional Requirements}
The project also has important non-functional constraints:
\begin{itemize}
  \item \textbf{Real-time structure:} sensing, control, display, networking,
  inference, and cloud publishing must coexist without unacceptable blocking.
  \item \textbf{Responsiveness:} web interactions should feel immediate to a
  human user, while periodic tasks should keep predictable cadence.
  \item \textbf{Modularity:} different subsystems should be readable and
  maintainable as separate source files and tasks.
  \item \textbf{Reduced shared global state:} the project should avoid broad
  global-variable usage and instead favor queues, semaphores, and protected
  shared state.
  \item \textbf{Fault tolerance:} the system should tolerate Wi-Fi disconnects,
  STA connection failure, and sensor read failure without total application
  collapse.
  \item \textbf{Scalability:} the architecture should be extendable with more
  sensors, more RPC methods, or richer dashboards in the future.
\end{itemize}

\subsection{Hardware and Software Environment}
The current implementation is designed around the following hardware and
software stack:
\begin{itemize}
  \item \textbf{Board:} ESP32-S3-based Yolo UNO platform.
  \item \textbf{Sensor:} DHT20 temperature and humidity sensor on the
  I\textsuperscript{2}C bus.
  \item \textbf{LCD:} 16x2 I\textsuperscript{2}C LCD at address
  \texttt{0x21}.
  \item \textbf{Discrete LED:} GPIO48 for temperature-dependent blink logic.
  \item \textbf{NeoPixel:} GPIO45, one RGB LED for humidity visualization.
  \item \textbf{External user-controlled LED:} GPIO5 with PWM control from the
  web interface.
  \item \textbf{Build environment:} PlatformIO with Arduino framework.
  \item \textbf{Software components:} FreeRTOS primitives, WebServer,
  Preferences, PubSubClient, ArduinoJson, and TensorFlow Lite Micro
  \cite{platformio,freertos_kernel,pubsubclient,arduinojson,tflm}.
\end{itemize}

These hardware and software choices define the practical limits of the system
and shape the architecture described in the next section.
